\documentclass[12pt,a4paper]{article}

% -- Packages ------------------------------------------------------------------
\usepackage[a4paper, top=2cm, bottom=2cm, left=2.2cm, right=2.2cm]{geometry}
\usepackage[utf8]{inputenc}
\usepackage[T1]{fontenc}
\usepackage{lmodern}
\usepackage{microtype}
\usepackage{xcolor}
\usepackage{listings}
\usepackage{mdframed}
\usepackage{enumitem}
\usepackage{titlesec}
\usepackage{titletoc}
\usepackage{fancyhdr}
\usepackage{booktabs}
\usepackage{array}
\usepackage{tabularx}
\usepackage{multicol}
\usepackage{amsmath}
\usepackage{amssymb}
\usepackage{graphicx}
\usepackage{hyperref}
\usepackage{tcolorbox}
\usepackage{fontawesome5}
\usepackage{parskip}
\usepackage{soul}
\usepackage{tikz}
\usepackage{pgffor}

\tcbuselibrary{skins, breakable, listingsutf8}

% -- Colours -------------------------------------------------------------------
\definecolor{pyblue}{HTML}{1F4E79}
\definecolor{pyteal}{HTML}{006B6B}
\definecolor{pyorange}{HTML}{C05000}
\definecolor{pygreen}{HTML}{1A6B1A}
\definecolor{pyred}{HTML}{8B0000}
\definecolor{pygray}{HTML}{F4F6F9}
\definecolor{pyborder}{HTML}{CCCCCC}
\definecolor{codebg}{HTML}{F0F2F5}
\definecolor{codeframe}{HTML}{C8D0DC}
\definecolor{kwcolor}{HTML}{0000CD}
\definecolor{strcolor}{HTML}{A31515}
\definecolor{cmtcolor}{HTML}{608B4E}
\definecolor{numcolor}{HTML}{098658}
\definecolor{builtincolor}{HTML}{795E26}
\definecolor{noteblue}{HTML}{E8F4FD}
\definecolor{noteborder}{HTML}{2196F3}
\definecolor{warnbg}{HTML}{FFF8E1}
\definecolor{warnborder}{HTML}{FF8F00}
\definecolor{tipbg}{HTML}{E8F5E9}
\definecolor{tipborder}{HTML}{388E3C}
\definecolor{dangerbg}{HTML}{FFEBEE}
\definecolor{dangerborder}{HTML}{C62828}

% -- Listings (Python syntax highlight) ----------------------------------------
\lstdefinestyle{python}{
  language=Python,
  backgroundcolor=\color{codebg},
  frame=single,
  framesep=6pt,
  rulecolor=\color{codeframe},
  basicstyle=\ttfamily\small,
  keywordstyle=\color{kwcolor}\bfseries,
  stringstyle=\color{strcolor},
  commentstyle=\color{cmtcolor}\itshape,
  numberstyle=\tiny\color{gray},
  numbers=left,
  numbersep=8pt,
  stepnumber=1,
  tabsize=4,
  showspaces=false,
  showstringspaces=false,
  breaklines=true,
  breakatwhitespace=false,
  captionpos=b,
  morekeywords={self, cls, None, True, False, lambda, yield, async, await},
  emph={print, range, len, type, input, int, float, str, list, dict, set, tuple, super, zip, map, filter, enumerate, sorted, reversed, open, isinstance},
  emphstyle=\color{builtincolor}\bfseries,
  xleftmargin=10pt,
  xrightmargin=4pt,
  aboveskip=8pt,
  belowskip=8pt,
}
\lstset{style=python}

% -- Inline code ---------------------------------------------------------------
\newcommand{\code}[1]{\texttt{\small\color{pyblue}#1}}

% -- tcolorbox environments ----------------------------------------------------
\tcbset{
  base/.style={
    breakable, enhanced, fonttitle=\bfseries\small,
    left=6pt, right=6pt, top=4pt, bottom=4pt,
    boxrule=0.5pt,
  }
}

\newtcolorbox{notebox}[1][]{
  base, colback=noteblue, colframe=noteborder,
  title={\faInfoCircle\ Note}, #1
}
\newtcolorbox{warnbox}[1][]{
  base, colback=warnbg, colframe=warnborder,
  title={\faExclamationTriangle\ Warning}, #1
}
\newtcolorbox{tipbox}[1][]{
  base, colback=tipbg, colframe=tipborder,
  title={\faLightbulb\ Tip}, #1
}
\newtcolorbox{dangerbox}[1][]{
  base, colback=dangerbg, colframe=dangerborder,
  title={\faBomb\ Watch Out}, #1
}
\newtcolorbox{conceptbox}[2][]{
  base, colback=pygray, colframe=pyblue,
  title={\faCode\ #2}, #1
}
\newtcolorbox{summarybox}[1][]{
  base, colback=pygray, colframe=pyteal,
  title={Section Summary}, #1
}

% -- Section formatting --------------------------------------------------------
\titleformat{\section}
  {\color{pyblue}\Large\bfseries}
  {\thesection.}{0.5em}{}
  [\color{pyblue}\titlerule]
\titleformat{\subsection}
  {\color{pyteal}\large\bfseries}
  {\thesubsection}{0.5em}{}
\titleformat{\subsubsection}
  {\color{pyorange}\normalsize\bfseries}
  {\thesubsubsection}{0.5em}{}

\titlespacing*{\section}{0pt}{18pt}{8pt}
\titlespacing*{\subsection}{0pt}{12pt}{4pt}
\titlespacing*{\subsubsection}{0pt}{8pt}{4pt}

% -- Header / Footer -----------------------------------------------------------
\pagestyle{fancy}
\fancyhf{}
\fancyhead[L]{\small\color{pyblue}\bfseries Python -- Lecture Notes}
\fancyhead[R]{\small\color{gray}\thesection\ \leftmark}
\fancyfoot[C]{\small\color{gray}\thepage}
\setlength{\headheight}{14pt}
\renewcommand{\headrulewidth}{0.5pt}
\renewcommand{\footrulewidth}{0.3pt}

% -- Hyperref ------------------------------------------------------------------
\hypersetup{
  colorlinks=true, linkcolor=pyblue, urlcolor=pyteal,
  pdftitle={Python Lecture Notes}, pdfauthor={Student Notes}
}

% -- List settings -------------------------------------------------------------
\setlist[itemize]{leftmargin=1.5em, topsep=2pt, itemsep=1pt, parsep=0pt}
\setlist[enumerate]{leftmargin=1.5em, topsep=2pt, itemsep=1pt, parsep=0pt}

% -----------------------------------------------------------------------------
\begin{document}
% -----------------------------------------------------------------------------

% -- Title page ----------------------------------------------------------------
\begin{titlepage}
  \centering
  \vspace*{2cm}
  {\Huge\bfseries\color{pyblue} Python -- Complete Lecture Notes \par}
  \vspace{0.5cm}
  {\large\color{gray} From Core Syntax to OOP, Functional Programming \& Concurrency \par}
  \vspace{1.5cm}
  \begin{tcolorbox}[colback=pyblue!8, colframe=pyblue, width=0.8\textwidth]
    \centering
    \textit{``Python is designed so that any developer can learn it.\\
    It is meant for everyone, not just specialists.''} \\[4pt]
    \small --- Guido van Rossum, Creator of Python (1991)
  \end{tcolorbox}
  \vspace{1.5cm}
  {\large\bfseries Topics Covered \par}
  \vspace{0.4cm}
  \begin{multicols}{2}
    \begin{itemize}[leftmargin=1em]
      \item Why Python? Language overview
      \item Data types \& memory model
      \item Operators \& expressions
      \item Strings \& collections
      \item Control flow \& loops
      \item Functions \& closures
      \item Modules \& packages
      \item OOP -- Classes \& objects
      \item Inheritance \& polymorphism
      \item Abstraction \& decorators
      \item Exception handling
      \item Threads \& multiprocessing
    \end{itemize}
  \end{multicols}
  \vfill
  {\small\color{gray} Personal study notes -- Intermediate level \par}
\end{titlepage}

% -- Table of Contents ---------------------------------------------------------
\tableofcontents
\newpage

% -----------------------------------------------------------------------------
\section{Python at a Glance}
% -----------------------------------------------------------------------------

\subsection{What is Python?}

Python is a \textbf{high-level, general-purpose, multi-paradigm} programming language first released in \textbf{1991} by Guido van Rossum. ``High-level'' means it reads close to English -- you don't write ones and zeros or assembly instructions. The interpreter converts your code to machine code at runtime.

\begin{notebox}
\textbf{Why ``high-level'' matters:}
Assembly languages were hardware-close but painful. C/C++ improved things but are complex. Python's goal was simplicity -- anyone should be able to program, not just specialists.
\end{notebox}

\subsection{Multi-Paradigm Support}

Python doesn't force one style. You can write code in whichever style fits your problem:

\begin{center}
\begin{tabular}{lll}
\toprule
\textbf{Paradigm} & \textbf{Core Idea} & \textbf{Examples in other languages} \\
\midrule
Procedural & Step-by-step instructions, functions & C \\
Object-Oriented & Model reality as objects & Java, C\# \\
Functional & Functions as values, avoid state & Haskell, JS \\
Structured & Blocks, conditions, loops & Most languages \\
\bottomrule
\end{tabular}
\end{center}

\subsection{Why Python is Famous}

\begin{itemize}
  \item \textbf{Easiest to read/write:} Syntax feels like English.
  \item \textbf{Huge community:} Millions of developers = tons of libraries.
  \item \textbf{AI/ML ecosystem:} NumPy, Pandas, TensorFlow, PyTorch -- all Python.
  \item \textbf{Web development:} Django, Flask, Tornado.
  \item \textbf{System scripting, automation, data science, GUIs.}
\end{itemize}

\subsection{How Python Runs}

\begin{enumerate}
  \item You write \code{.py} files.
  \item Python's \textbf{CPython interpreter} compiles to bytecode (\code{.pyc}).
  \item The Python Virtual Machine (PVM) executes bytecode.
  \item OS and hardware execute the final machine instructions.
\end{enumerate}

\begin{tipbox}
You can test code in the Python REPL (interactive prompt) by typing \code{python} in your terminal. The \texttt{>>>} prompt means you're in REPL mode. Type \code{exit()} to leave.
\end{tipbox}

% -----------------------------------------------------------------------------
\section{Variables, Memory Model \& the Object System}
% -----------------------------------------------------------------------------

\subsection{Variables are References, Not Boxes}

This is the single most important mental model shift from C/Java. In Python, a variable is not a container holding a value. It is a \textbf{label} (reference) that points to an \textbf{object} in memory.

\begin{lstlisting}
a = 5        # Creates int object 5 in memory. 'a' points to it.
b = 5        # Python finds existing 5 object. 'b' ALSO points to same object.
print(id(a) == id(b))   # True -- same identity!

b = 6        # New int object 6 created. 'b' now points to 6.
print(id(a) == id(b))   # False -- different objects now.
\end{lstlisting}

\begin{notebox}
\textbf{id()} returns a unique integer identifying an object's memory location. Two objects with the same id are the same object.
\end{notebox}

\subsection{Garbage Collection}

When no variable points to an object, Python's garbage collector reclaims its memory. You don't manage memory manually.

\begin{lstlisting}
a = 5    # object 5, referenced by a
b = 5    # same object 5, referenced by a and b
b = 6    # object 6 for b. Object 5 still has 'a'.
a = 10   # Now object 5 has ZERO references -> garbage collected.
\end{lstlisting}

\subsection{String Interning}

Python caches small integers (typically \(-5\) to \(256\)) and short strings as an optimization. This means \code{a = 5; b = 5} share the same object. For large strings or numbers, Python creates separate objects even with identical values.

\begin{lstlisting}
a = "hello"
b = "hello"
print(id(a) == id(b))   # Usually True -- interned (short, simple string)

a = "my favorite color is blue"
b = "my favorite color is blue"
print(id(a) == id(b))   # False -- too long, not interned
\end{lstlisting}

\subsection{Multiple Assignment}

Python supports elegant multiple assignment using tuple packing/unpacking:

\begin{lstlisting}
x, y = 6, 5          # Assign in one line
x, y = y, x          # Swap! Behind the scenes: tuple pack then unpack.
a = b = c = 0        # Chain assignment -- all point to same 0 object
\end{lstlisting}

% -----------------------------------------------------------------------------
\section{Data Types}
% -----------------------------------------------------------------------------

\subsection{Numeric Types}

\begin{lstlisting}
x = 10          # int -- unlimited precision in Python
pi = 3.14       # float -- IEEE 754 double precision
c = 3 + 4j      # complex -- uses 'j' not 'i'

# Type conversion
k = int(pi)     # 3  -- truncates, doesn't round
p = float(k)    # 3.0
c = complex(2, 3)  # (2+3j)

# Check
print(type(x))  # <class 'int'>
\end{lstlisting}

\begin{notebox}
\textbf{Boolean is a subtype of int:}
\code{True == 1} and \code{False == 0}. So \code{int(True)} gives 1. You can even do \code{True + True == 2}.
\end{notebox}

\subsection{Strings}

Strings are \textbf{immutable sequences of characters}. Once created, characters cannot be changed in-place.

\begin{lstlisting}
name = 'Naven'          # single quotes
msg  = "Hello world"    # double quotes
para = """Line one
Line two"""             # triple quotes -- multi-line

# Indexing (0-based)
print(name[0])     # 'N'   (first character)
print(name[-1])    # 'n'   (last character)

# Slicing [start:stop:step]  -- stop is EXCLUSIVE
text = 'telusko'
print(text[3:7])   # 'usko'  -- indices 3,4,5,6
print(text[:3])    # 'tel'   -- from start to index 2
print(text[3:])    # 'usko'  -- from index 3 to end
print(text[-4:])   # 'usko'  -- last 4 characters
print(text[::2])   # 'tlso'  -- every second character

# Repetition
print('Na' * 3)    # 'NaNaNa'

# f-strings (preferred formatting)
age = 30
print(f"My name is {name} and I am {age} years old.")
print(f"Pi is approximately {3.14159:.2f}")  # 2 decimal places
\end{lstlisting}

\subsubsection{Escape Sequences}

\begin{center}
\begin{tabular}{ll}
\toprule
\textbf{Escape} & \textbf{Meaning} \\
\midrule
\texttt{\textbackslash n} & Newline \\
\texttt{\textbackslash t} & Tab \\
\texttt{\textbackslash\textbackslash} & Literal backslash \\
\texttt{\textbackslash '} & Single quote inside single-quoted string \\
\texttt{\textbackslash "} & Double quote inside double-quoted string \\
\bottomrule
\end{tabular}
\end{center}

\begin{lstlisting}
# Problem: Windows path
path = "C:\\Users\\Naven"   # double backslash to escape
# OR use raw string:
path = r"C:\Users\Naven"    # r prefix means no escape processing
\end{lstlisting}

\subsubsection{Key String Methods}

\begin{lstlisting}
s = "  Hello, Python!  "
print(s.strip())         # "Hello, Python!"  -- remove whitespace
print(s.lower())         # "  hello, python!  "
print(s.upper())         # "  HELLO, PYTHON!  "
print(s.replace("Python", "World"))  # "  Hello, World!  "
print(s.split(","))      # ["  Hello", " Python!  "]
print(",".join(["a", "b", "c"]))  # "a,b,c"
print(s.find("Python"))  # index of first match (-1 if not found)
print(len(s))            # length including spaces
\end{lstlisting}

\subsection{Lists}

Lists are \textbf{ordered, mutable, allow duplicates} and can hold mixed types.

\begin{lstlisting}
nums = [45, 87, 21, 24, 99]      # list of integers
mix  = [1, "hello", 3.14, True]  # mixed types -- totally fine

# Access
print(nums[0])     # 45
print(nums[-1])    # 99

# Slicing works same as strings
print(nums[1:3])   # [87, 21]

# Common methods
nums.append(33)    # Add at end
nums.insert(1, 55) # Insert 55 at index 1; shifts others right
nums.remove(87)    # Remove first occurrence of value 87
val = nums.pop()   # Remove and return LAST element
val = nums.pop(2)  # Remove and return element at index 2
nums.reverse()     # Reverse in-place
nums.sort()        # Sort in-place (ascending by default)
count = nums.count(21)   # Number of times 21 appears
idx   = nums.index(24)   # Index of first occurrence of 24

# del keyword for slices
del nums[2:4]      # Remove elements at index 2 and 3

# Extend vs Append
nums.extend([100, 200])  # Adds each element from iterable
nums.append([100, 200])  # Adds the LIST as a single element

# Python built-in functions on lists
print(min(nums), max(nums), sum(nums), len(nums))
\end{lstlisting}

\begin{notebox}
\textbf{List inside list (nested):} \code{mix = [nums, names]} -- each inner list is a single element. Access with \code{mix[0][2]} -- first list, third element.
\end{notebox}

\subsection{Tuples}

Tuples are \textbf{ordered and immutable}. Use them when data should not change.

\begin{lstlisting}
t = (23, 45, 67, 43)   # with parentheses
t = 23, 45, 67, 43     # parentheses optional -- same result

# Single-element tuple NEEDS trailing comma
single = (42,)   # This is a tuple
not_tuple = (42) # This is just int 42 in parentheses

# Access and slicing work same as list
print(t[1])     # 45
print(t[-1])    # 43
print(t[1:3])   # (45, 67)

# Only two methods (can't mutate)
print(t.count(45))  # occurrences
print(t.index(67))  # index of value

# Immutable = cannot reassign elements
# t[0] = 99  -> TypeError: tuple does not support item assignment

# Tuple UNPACKING
num, name, score = (1, "Naven", 95.5)
print(num, name, score)

# Multiple return values use tuple implicitly
def minmax(lst):
    return min(lst), max(lst)   # returns a tuple

lo, hi = minmax([3, 1, 7, 2])
\end{lstlisting}

\begin{tipbox}
\textbf{When to use tuple vs list?} Tuple = fixed data (coordinates, RGB colour, database row). List = collection that grows/shrinks. Tuples are slightly faster and signal ``this won't change''.
\end{tipbox}

\subsection{Sets}

Sets are \textbf{unordered collections of unique elements}. Perfect for membership testing and eliminating duplicates.

\begin{lstlisting}
s1 = {23, 56, 78, 21, 56}  # duplicate 56 silently removed
print(s1)   # {56, 21, 78, 23} -- order may vary (hashed internally)

# Empty set MUST use set() -- {} creates an empty DICT
empty = set()

# Membership test -- O(1) speed due to hashing
print(21 in s1)   # True
print(99 in s1)   # False

# Set operations (mirrors math set theory)
s2 = {'a', 'b', 'c', 'd', 'm'}
s3 = {'a', 'e', 'i', 'o', 'u', 'd', 'p'}

print(s2 - s3)    # Difference: in s2 but NOT in s3
print(s2 | s3)    # Union: all elements from both
print(s2 & s3)    # Intersection: only common elements
print(s2 ^ s3)    # Symmetric difference: in one but NOT both
\end{lstlisting}

\begin{notebox}
\textbf{Why unordered?} Each element is hashed. Python uses the hash to decide where to store it -- this makes lookup O(1) instead of O(n). The side effect: order is not preserved.
\end{notebox}

\subsection{Dictionaries}

Dictionaries are \textbf{key-value pairs}. Keys form a set (unique, unordered). Values form a list (can repeat). Think of a real dictionary: word -> meaning.

\begin{lstlisting}
phone_book = {
    "Naven":  9876543210,
    "Kiran":  8765432109,
    "Hush":   7654321098,
}

# Access by key (KeyError if missing)
print(phone_book["Naven"])      # 9876543210

# .get() -- safe access, returns None (or default) if key missing
print(phone_book.get("Naven"))         # 9876543210
print(phone_book.get("Unknown"))       # None
print(phone_book.get("Unknown", -1))   # -1 (custom default)

# Add / update
phone_book["Social"] = 6543210987  # Add new key
phone_book["Naven"] = 1111111111   # Update existing key

# Delete
phone_book.pop("Hush")       # Remove and return value
del phone_book["Social"]     # Remove without returning

# Iteration
for key in phone_book:
    print(key, "->", phone_book[key])

for key, val in phone_book.items():   # Pythonic way
    print(f"{key}: {val}")

# Nested dict / dict with list values
data = {
    "js":     "VSCode",
    "python": ["VSCode", "PyCharm"],
    "java":   {"core": "VSCode", "spring": "IntelliJ"}
}
print(data["java"]["spring"])   # "IntelliJ"
\end{lstlisting}

\subsubsection{Creating from Sets + Lists}

\begin{lstlisting}
keys   = {"Naven", "Hitesh", "Shamik"}  # set -- unique keys
values = [34, 64, 87]                   # list of values

# zip + dict -- combines into dict
d = dict(zip(keys, values))
# Note: set order is not guaranteed, so key-value pairing may vary.
# Use a list for keys if order matters:
keys = ["Naven", "Hitesh", "Shamik"]
d = dict(zip(keys, values))
print(d)   # {'Naven': 34, 'Hitesh': 64, 'Shamik': 87}
\end{lstlisting}

% -----------------------------------------------------------------------------
\section{Operators}
% -----------------------------------------------------------------------------

\subsection{Arithmetic Operators}

\begin{center}
\begin{tabular}{lll}
\toprule
\textbf{Operator} & \textbf{Name} & \textbf{Example} \\
\midrule
\texttt{+} & Addition & \texttt{3 + 5 = 8} \\
\texttt{-} & Subtraction & \texttt{5 - 2 = 3} \\
\texttt{*} & Multiplication & \texttt{3 * 4 = 12} \\
\texttt{/} & True Division & \texttt{8 / 2 = 4.0} (always float) \\
\texttt{//} & Floor Division & \texttt{9 // 2 = 4} (quotient) \\
\texttt{\%} & Modulus & \texttt{9 \% 2 = 1} (remainder) \\
\texttt{**} & Exponentiation & \texttt{2 ** 10 = 1024} \\
\bottomrule
\end{tabular}
\end{center}

\begin{notebox}
Python follows \textbf{BODMAS/PEMDAS}. Multiplication happens before addition. Use parentheses to override: \code{(9 + 2) * 3 = 33} vs \code{9 + 2 * 3 = 15}.
\end{notebox}

\subsection{Comparison (Relational) Operators}

These always return a \code{bool} (\code{True} or \code{False}).

\begin{lstlisting}
a, b = 4, 5
print(a > b)    # False
print(a < b)    # True
print(a >= b)   # False
print(a <= b)   # False
print(a == b)   # False  -- double equals for comparison
print(a != b)   # True
\end{lstlisting}

\begin{dangerbox}
\code{=} is \textit{assignment}. \code{==} is \textit{comparison}. Confusing them is the most common beginner error.
\end{dangerbox}

\subsection{Logical Operators -- Truth Table}

\begin{center}
\begin{tabular}{ccccc}
\toprule
\textbf{C1} & \textbf{C2} & \textbf{C1 and C2} & \textbf{C1 or C2} & \textbf{not C1} \\
\midrule
True & True & True & True & False \\
True & False & \textbf{False} & True & False \\
False & True & \textbf{False} & True & True \\
False & False & \textbf{False} & \textbf{False} & True \\
\bottomrule
\end{tabular}
\end{center}

\begin{lstlisting}
a, b = 4, 5
print(a < 10 and b > 1)    # True AND True  -> True
print(a < 10 and b > 10)   # True AND False -> False (and needs both)
print(a < 10 or  b > 10)   # True OR  False -> True  (or needs just one)
print(not (a == b))         # not False -> True
\end{lstlisting}

\subsection{Assignment Shorthand}

\begin{lstlisting}
a = 5
a += 2   # a = a + 2  ->  7
a -= 1   # a = a - 1  ->  6
a *= 3   # a = a * 3  ->  18
a //= 4  # a = a // 4 ->  4
a **= 2  # a = a ** 2 ->  16
a %= 5   # a = a % 5  ->  1
\end{lstlisting}

% -----------------------------------------------------------------------------
\section{Control Flow}
% -----------------------------------------------------------------------------

\subsection{if / elif / else}

Python uses \textbf{indentation} (not curly braces) to define blocks. This enforces readable code.

\begin{lstlisting}
num = int(input("Enter a number: "))

if num % 2 == 0:
    print("Even")
elif num % 2 != 0:
    print("Odd")
else:
    print("This won't run")   # redundant here but shows else syntax
\end{lstlisting}

\begin{tipbox}
The \code{else} in an if/elif chain is optional. \code{elif} can appear as many times as needed. There can only be one \code{else}. As soon as a condition is \code{True}, the rest are skipped.
\end{tipbox}

\subsubsection{Nested if}

\begin{lstlisting}
num = int(input("Enter a number: "))
salary = int(input("Enter salary: "))

if num % 2 == 0:
    print("Even number")
    if salary >= 5:
        print("Great job!")
    else:
        print("Better luck next time.")
else:
    print("Odd number")
\end{lstlisting}

\subsection{match Statement (Python 3.10+)}

A cleaner alternative to long \code{elif} chains when matching specific values:

\begin{lstlisting}
num = 3

match num:
    case 1: print("One")
    case 2: print("Two")
    case 3: print("Three")
    case 4: print("Four")
    case 5: print("Five")
    case _: print("Number out of range")   # _ is the default/wildcard
\end{lstlisting}

\begin{notebox}
Unlike C's \texttt{switch}, Python's \texttt{match} does NOT fall through. No \code{break} needed. The \code{\_} case is the default -- like \code{else}.
\end{notebox}

\subsection{while Loop}

Repeats a block as long as a condition is \code{True}.

\begin{lstlisting}
i = 1
while i <= 5:
    print(f"Hello world {i}")
    i += 1           # CRITICAL: without this, infinite loop!
print(f"Final value of i: {i}")   # 6  (exits when i=6 fails the check)
\end{lstlisting}

\subsubsection{Nested while -- ``Hour and Day'' Analogy}

\begin{lstlisting}
i = 1
while i <= 3:          # outer loop (days)
    j = 1              # reset j for every outer iteration (day resets hours)
    while j <= 4:      # inner loop (hours)
        print("Telusko " * 1, end="")
        print("rocks " * 1, end="")
        j += 1
    print()            # newline after each outer iteration
    i += 1
\end{lstlisting}

\begin{dangerbox}
\textbf{Infinite loop trap:} If you forget to initialize \code{j} inside the outer loop, after the first outer iteration \code{j} stays at 5, and the inner loop never executes again.
\end{dangerbox}

\subsection{for Loop}

The \code{for} loop iterates over any \textbf{iterable} (list, tuple, string, range, dict, etc.).

\begin{lstlisting}
# Over a list
data = [2, "Naven", 4.5, 8, "Telusko", "Python"]
for value in data:
    print(value)

# Over a range
for i in range(10):          # 0, 1, 2, ..., 9
    print(i)

for i in range(2, 11, 2):   # even numbers: 2, 4, 6, 8, 10
    print(i)

# Over dict (keys by default)
for key in phone_book:
    print(key, "->", phone_book[key])

# Over string (character by character)
for ch in "Python":
    print(ch)
\end{lstlisting}

\subsection{break and continue}

\begin{lstlisting}
# continue -- skip current iteration
for i in range(10):
    if i % 3 == 0:
        continue       # skip multiples of 3
    print(i)           # prints 1, 2, 4, 5, 7, 8

# break -- exit loop entirely
for i in range(10):
    if i == 5:
        break          # stop at 5
    print(i)           # prints 0, 1, 2, 3, 4

print("bye")           # always prints (outside loop)
\end{lstlisting}

\begin{tipbox}
\textbf{for...else and while...else:} Python has an optional \code{else} clause on loops. The \code{else} block runs if the loop completed \textit{without} hitting a \code{break}.
\end{tipbox}

% -----------------------------------------------------------------------------
\section{Functions}
% -----------------------------------------------------------------------------

\subsection{Why Functions?}

Functions let you \textbf{write code once, call it many times}. They also give a name to a block of logic, making code readable. Python's built-ins like \code{print()}, \code{len()}, \code{int()} are all functions defined elsewhere.

\begin{lstlisting}
# Without function -- repetitive
a, b = 4, 5
c = a + b
print(c)
# ... 50 lines later, same logic again
a, b = 10, 20
c = a + b
print(c)

# With function -- DRY (Don't Repeat Yourself)
def add(x, y):
    """Adds two values and returns the result."""
    return x + y

result = add(4, 5)
print(result)        # 9
result = add(10, 20)
print(result)        # 30
\end{lstlisting}

\begin{notebox}
The triple-quoted string immediately after \code{def} is a \textbf{docstring}. It documents what the function does. Access it with \code{add.\_\_doc\_\_} or hover in your IDE.
\end{notebox}

\subsection{Parameters and Arguments}

\subsubsection{Default Arguments}

\begin{lstlisting}
def greet(name, greeting="Hello"):
    print(f"{greeting}, {name}!")

greet("Naven")              # Hello, Naven!   (uses default)
greet("Naven", "Namaste")   # Namaste, Naven!  (overrides default)
\end{lstlisting}

\subsubsection{Keyword Arguments}

When calling a function, you can specify arguments by name -- order doesn't matter.

\begin{lstlisting}
def person(name, age=18):
    print(f"Name: {name}, Age: {age}")

person(age=30, name="Naven")  # keyword args -- order flexible
\end{lstlisting}

\subsubsection{Variable-Length Positional Arguments (*args)}

\begin{lstlisting}
def add(first, *rest):
    """first is a regular arg. rest is a tuple of extra args."""
    total = first
    for n in rest:
        total += n
    return total

print(add(4, 5, 6, 7))   # 22  -- works for any number of extra args
\end{lstlisting}

\subsubsection{Variable-Length Keyword Arguments (**kwargs)}

\begin{lstlisting}
def person(name, **extra_info):
    """extra_info is a dict of any keyword args passed."""
    print(f"Name: {name}")
    for key, val in extra_info.items():
        print(f"  {key}: {val}")

person("Naven", age=30, location="Pune", tech="Python")
# Name: Naven
#   age: 30
#   location: Pune
#   tech: Python
\end{lstlisting}

\subsection{Return Values and None}

\begin{lstlisting}
def add(a, b):
    return a + b          # explicit return value

def greet(name):
    print(f"Hello {name}")  # no return -- implicitly returns None

result = greet("Naven")
print(result)              # None

# Functions can return multiple values (as a tuple)
def divmod_manual(a, b):
    return a // b, a % b   # quotient and remainder

q, r = divmod_manual(9, 2)
print(q, r)   # 4  1
\end{lstlisting}

\subsection{Scope: Global vs Local Variables}

\begin{lstlisting}
a = 10   # global variable

def something():
    a = 15   # This creates a NEW local variable 'a' -- does NOT change global!
    print(f"Inside function: {a}")   # 15

something()
print(f"Outside function: {a}")   # 10 -- global unchanged

# To modify the global variable inside a function:
def change_global():
    global a
    a = 20   # Now this modifies the global 'a'

change_global()
print(a)   # 20

# globals() -- returns a dict of all current global variables
print(globals()['a'])   # 20
\end{lstlisting}

\begin{warnbox}
Using \code{global} is generally considered bad practice in large codebases. Prefer passing values as arguments and returning them. Reserve \code{global} for truly necessary cases.
\end{warnbox}

% -----------------------------------------------------------------------------
\section{Functional Programming Tools}
% -----------------------------------------------------------------------------

\subsection{Functions as First-Class Objects}

In Python, functions are objects. They can be:
\begin{itemize}
  \item Assigned to variables
  \item Passed as arguments to other functions
  \item Returned from functions
  \item Stored in lists/dicts
\end{itemize}

This is the foundation of functional programming in Python.

\subsection{Lambda (Anonymous Functions)}

A \code{lambda} is a one-liner function -- no name, no \code{def}, no \code{return} statement needed.

\begin{lstlisting}
# Regular function
def square(n): return n * n

# Lambda equivalent
square = lambda n: n * n

# Multi-parameter lambda
add = lambda a, b: a + b
print(add(3, 4))   # 7

# Used directly in expressions
print((lambda x: x ** 2)(5))   # 25
\end{lstlisting}

\begin{tipbox}
Use lambdas for short, one-off functions, especially inside \code{map()}, \code{filter()}, \code{reduce()}, or when passing a function to another function.
\end{tipbox}

\subsection{map, filter, reduce}

\begin{lstlisting}
from functools import reduce

nums = [4, 2, 9, 7, 6, 8, 3, 1]

# filter -- keep elements where function returns True
evens = list(filter(lambda n: n % 2 == 0, nums))
print(evens)   # [4, 2, 6, 8]

# map -- transform each element
doubles = list(map(lambda n: n * 2, evens))
print(doubles)   # [8, 4, 12, 16]

# reduce -- fold all elements into one value
total = reduce(lambda a, b: a + b, doubles)
print(total)   # 40

# Combining in one pipeline:
result = reduce(
    lambda a, b: a + b,
    map(lambda n: n * 2,
        filter(lambda n: n % 2 == 0, nums))
)
print(result)   # 40
\end{lstlisting}

\subsection{Higher-Order Functions}

A function that \textbf{takes or returns another function} is called a higher-order function.

\begin{lstlisting}
def square(n):  return n * n
def cube(n):    return n * n * n

def operate(nums, operation):
    """Higher-order function -- accepts any single-arg function."""
    for n in nums:
        result = operation(n)
        print(f"{operation.__name__}({n}) = {result}")

operate([5, 6, 7], square)
operate([5, 6, 7], cube)
operate([5, 6, 7], lambda n: n + 100)   # anonymous inline
\end{lstlisting}

\subsection{Inner Functions and Closures}

\begin{lstlisting}
def outer():
    def inner(x):
        print(f"Inner says: {x}")
    inner(42)          # can call inner from within outer
    return inner       # return the function OBJECT (not call it)

fn = outer()   # fn is now the inner function
fn(99)         # calls inner(99) -- Inner says: 99

# Closure -- inner function captures outer scope
def counter_factory(start):
    count = start
    def increment():
        nonlocal count   # nonlocal lets inner modify outer's var
        count += 1
        return count
    return increment

counter = counter_factory(10)
print(counter())   # 11
print(counter())   # 12
print(counter())   # 13
\end{lstlisting}

\subsection{Decorators}

A decorator is a function that \textbf{wraps another function} to add behavior before/after it, without modifying the original function.

\begin{lstlisting}
def log_deco(func):
    """Decorator that logs function calls."""
    def wrap(*args, **kwargs):
        print(f"Calling '{func.__name__}' with args={args}")
        result = func(*args, **kwargs)
        print(f"Result: {result}")
        return result
    return wrap

@log_deco          # Syntactic sugar: divide = log_deco(divide)
def divide(a, b):
    return a / b

divide(10, 2)
# Calling 'divide' with args=(10, 2)
# Result: 5.0
\end{lstlisting}

\begin{lstlisting}
# Decorator that enforces "greater number first" behavior
def greater_first(func):
    def wrap(a, b):
        if a < b:
            a, b = b, a     # swap so larger is always first
        return func(a, b)
    return wrap

@greater_first
def divide(a, b):
    return a / b

@greater_first
def subtract(a, b):
    return a - b

print(divide(2, 8))    # 8/2 = 4.0  (not 2/8!)
print(subtract(3, 9))  # 9-3 = 6
\end{lstlisting}

\begin{notebox}
\textbf{Why decorators matter in real projects:} Frameworks like Flask/Django use decorators heavily -- \code{@app.route('/home')} tells Flask which function handles the \code{/home} URL. \code{@login\_required} protects routes. You write the logic; the decorator adds infrastructure.
\end{notebox}

% -----------------------------------------------------------------------------
\section{Modules, Packages \& \texttt{\_\_name\_\_}}
% -----------------------------------------------------------------------------

\subsection{Modules}

A \textbf{module} is any \code{.py} file. When it contains only function/class definitions (no top-level execution), it acts as a library for other files.

\begin{lstlisting}
# calc.py (a module)
def add(a, b):
    return a + b

def subtract(a, b):
    return a - b
\end{lstlisting}

\begin{lstlisting}
# demo.py (uses the module)
import calc                      # import whole module
print(calc.add(3, 4))            # 7 -- must prefix with module name

from calc import add             # import specific function
print(add(3, 4))                 # 7 -- no prefix needed

from calc import add, subtract   # import multiple
from calc import *               # import everything (avoid -- pollutes namespace)
import calc as c                 # alias
print(c.add(3, 4))
\end{lstlisting}

\subsection{Packages}

A \textbf{package} is a folder of modules. Group related modules into packages.

\begin{verbatim}
my_project/
  calculator/          <- package (folder)
    calc.py            <- module
    demo.py            <- module
  main.py              <- entry point
\end{verbatim}

\begin{lstlisting}
# main.py -- accessing from outside the package
from calculator import calc     # or
from calculator.calc import add
\end{lstlisting}

\subsection{The \texttt{\_\_name\_\_} Variable}

\begin{lstlisting}
# When you RUN a file directly:    __name__ == '__main__'
# When a file is IMPORTED:         __name__ == 'module_name'

# calc.py
def add(a, b):
    return a + b

if __name__ == '__main__':
    # This block only runs when calc.py is run directly
    # NOT when calc.py is imported by demo.py
    print(add(2, 3))   # test code
\end{lstlisting}

\begin{tipbox}
Always put test/demo code inside \code{if \_\_name\_\_ == '\_\_main\_\_':}. This way, your file works both as a module (importable) and as a standalone script.
\end{tipbox}

\subsection{The \texttt{math} Module -- Example}

\begin{lstlisting}
import math

print(math.sqrt(25))     # 5.0
print(math.ceil(59.4))   # 60  (ceiling -- always rounds UP)
print(math.floor(59.9))  # 59  (floor  -- always rounds DOWN)
print(math.pow(5, 3))    # 125.0 (5 to the power 3)
print(math.pi)           # 3.141592653589793

# help(math) or help(math.sqrt) shows docstrings in REPL
\end{lstlisting}

% -----------------------------------------------------------------------------
\section{Object-Oriented Programming (OOP)}
% -----------------------------------------------------------------------------

\subsection{The OOP Mindset}

OOP models software around \textbf{objects} rather than functions. Every real-world thing has:
\begin{itemize}
  \item \textbf{Properties} (what it \textit{knows}) -- stored as variables
  \item \textbf{Behaviour} (what it \textit{does}) -- stored as methods (functions inside class)
\end{itemize}

\textbf{Class = Blueprint.} \textbf{Object = Instance} created from that blueprint.

One blueprint -> many objects, just as one building plan -> many identical buildings.

\subsection{Defining a Class and Creating Objects}

\begin{lstlisting}
class Computer:
    brand = "Telusko"       # CLASS VARIABLE -- shared by all instances

    def __init__(self, cpu, ram, ssd):
        # INSTANCE VARIABLES -- unique per object
        self.cpu = cpu
        self.ram = ram
        self.ssd = ssd

    def config(self):
        """Print configuration of this computer."""
        print(f"Config: {self.cpu}, {self.ram} RAM, {self.ssd} SSD")

# Create objects (instances)
comp1 = Computer("i5",  "16GB", "512GB")
comp2 = Computer("i9",  "96GB", "2TB")

comp1.config()   # Config: i5,  16GB RAM, 512GB SSD
comp2.config()   # Config: i9,  96GB RAM,   2TB SSD

# Class variable accessed via class name or instance
print(Computer.brand)   # Telusko
print(comp1.brand)      # Telusko
\end{lstlisting}

\subsubsection{What is \texttt{self}?}

\code{self} is the first parameter of every instance method. It is a reference to \textit{the object calling the method}. Python passes it automatically -- you never pass \code{self} explicitly when calling.

\begin{lstlisting}
comp1.config()
# Python internally translates this to:
Computer.config(comp1)   # comp1 becomes 'self' inside the method
\end{lstlisting}

You can name it anything, but \code{self} is universally agreed upon. Use it.

\subsubsection{Three Types of Methods}

\begin{lstlisting}
class Computer:
    brand = "Telusko"

    def __init__(self, cpu, ram):
        self.cpu = cpu
        self.ram = ram

    # 1. Instance method -- works with instance variables via self
    def config(self):
        print(f"{self.cpu}, {self.ram}")

    # 2. Class method -- works with class variables via cls
    @classmethod
    def info(cls):
        print(f"Brand: {cls.brand}")

    # 3. Static method -- utility method, no self or cls
    @staticmethod
    def gb_to_bytes(gb):
        return gb * (1024 ** 3)

c = Computer("i7", "32GB")
c.config()               # instance method
Computer.info()          # class method
Computer.gb_to_bytes(16) # static method
\end{lstlisting}

\subsubsection{Constructor vs \texttt{\_\_init\_\_}}

\begin{lstlisting}
# __new__ is the true constructor -- creates the object
# __init__ is the initializer -- sets up the object's state
# Python calls __new__ then __init__ for every `ClassName()` call.
# You almost never need to define __new__ unless implementing singletons.

class A:
    def __new__(cls):
        print("Constructor called -- object being created")
        return super().__new__(cls)   # must return the object

    def __init__(self):
        print("Initializer called -- setting up state")

a = A()
# Constructor called -- object being created
# Initializer called -- setting up state
\end{lstlisting}

\subsection{Inheritance}

Inheritance lets one class \textbf{reuse} the variables and methods of another.

\begin{lstlisting}
class A:
    def feature1(self):  print("F1 works")
    def feature2(self):  print("F2 works")

class B(A):         # B inherits from A -- single inheritance
    def feature3(self):  print("F3 works")
    def feature4(self):  print("F4 works")

class C(B):         # C inherits from B which inherits from A -- multi-level
    def feature5(self):  print("F5 works")

obj = C()
obj.feature1()  # C -> B -> A -- found in A
obj.feature5()  # found in C itself
\end{lstlisting}

\subsubsection{Multiple Inheritance}

\begin{lstlisting}
class A:
    def show(self): print("in A show")

class B:
    def show(self): print("in B show")

class C(A, B):    # inherits from both A and B
    pass

c = C()
c.show()   # "in A show" -- MRO: C -> A -> B (left to right)

# To explicitly call B's show:
B.show(c)   # "in B show"
\end{lstlisting}

\subsubsection{Method Resolution Order (MRO)}

\begin{lstlisting}
print(C.__mro__)
# (<class 'C'>, <class 'A'>, <class 'B'>, <class 'object'>)
# Python searches left to right, then 'object' (root of all classes)
\end{lstlisting}

\subsubsection{\texttt{super()} and \texttt{\_\_init\_\_} Chains}

\begin{lstlisting}
class A:
    def __init__(self):
        print("A init")
    def feature1(self): print("F1 works")

class B(A):
    def __init__(self):
        super().__init__()   # call A's __init__
        print("B init")
    def feature3(self): print("F3 works")

obj = B()
# A init
# B init

obj.feature1()  # from A (inherited)
obj.feature3()  # from B
\end{lstlisting}

\subsection{Polymorphism}

``One interface, many forms.'' The same operation behaves differently depending on the object.

\subsubsection{Duck Typing}

\begin{lstlisting}
class Laptop:
    def build(self):
        print("Laptop builds")

class Desktop:
    def build(self):
        print("Desktop builds")

class Tablet:
    def open_pdf(self):       # NO build() method!
        print("Tablet opens PDF")

def developer_works(machine):
    machine.build()           # just needs a .build() method

developer_works(Laptop())     # works
developer_works(Desktop())    # works -- different class, same interface
developer_works(Tablet())     # AttributeError -- no .build()!
\end{lstlisting}

\begin{notebox}
\textbf{Duck typing:} ``If it quacks like a duck, it is a duck.'' Python doesn't check the type -- it just calls the method. If the method exists, it works.
\end{notebox}

\subsubsection{Operator Overloading}

You can define what \code{+}, \code{>}, \code{str()} etc. mean for your custom class by implementing dunder (magic) methods.

\begin{lstlisting}
class Account:
    def __init__(self, name, balance):
        self.name    = name
        self.balance = balance

    def __str__(self):
        """Called by print() / str()."""
        return f"{self.name}: {self.balance}"

    def __add__(self, other):
        """Defines what '+' means for two Account objects."""
        return Account("Combined", self.balance + other.balance)

    def __gt__(self, other):
        """Defines what '>' means."""
        return self.balance > other.balance

u1 = Account("Naven", 1000)
u2 = Account("Kiran", 2000)

print(u1)            # Naven: 1000  (uses __str__)
combined = u1 + u2   # uses __add__
print(combined)      # Combined: 3000
print(u1 > u2)       # False  (1000 > 2000)  (uses __gt__)
\end{lstlisting}

\subsubsection{Method Overriding}

When a child class redefines a parent method with the same name, the child's version is used.

\begin{lstlisting}
class A:
    def show(self): print("in A show")

class B(A):
    def show(self): print("in B show")   # OVERRIDES A's show

b = B()
b.show()   # "in B show" -- child's version wins
\end{lstlisting}

\subsection{Abstraction}

Abstraction means \textbf{hiding implementation details} and exposing only what's necessary. In Python, this is achieved through abstract classes (ABC).

\begin{lstlisting}
from abc import ABC, abstractmethod

class PaymentGateway(ABC):
    """Interface -- defines what every gateway MUST have."""

    @abstractmethod
    def pay(self, amount):
        pass   # no implementation -- just the contract

class RazorPay(PaymentGateway):
    def pay(self, amount):
        print(f"Paying Rs. {amount} via RazorPay")

class Telusko Pay(PaymentGateway):
    def pay(self, amount):
        print(f"Paying Rs. {amount} via TeluskoP ay")

# Can NOT instantiate abstract class:
# gw = PaymentGateway()  -> TypeError: Can't instantiate abstract class

rp = RazorPay()
rp.pay(500)   # Paying Rs. 500 via RazorPay
\end{lstlisting}

\begin{tipbox}
\textbf{Real-world use:} Your \code{checkout} page depends on \code{PaymentGateway}, not \code{RazorPay} specifically. Switch to a different gateway by just passing a different object -- no changes to checkout code. This is \textbf{loosely coupled design}.
\end{tipbox}

% -----------------------------------------------------------------------------
\section{Recursion}
% -----------------------------------------------------------------------------

\subsection{What is Recursion?}

A function calling itself. Every recursion needs:
\begin{enumerate}
  \item A \textbf{base case} -- the stopping condition (or it runs forever)
  \item A \textbf{recursive case} -- reducing the problem toward the base case
\end{enumerate}

\begin{lstlisting}
# Factorial: n! = n * (n-1) * ... * 1
def factorial(n):
    if n == 1:           # base case
        return 1
    return n * factorial(n - 1)   # recursive call

print(factorial(5))  # 120
# Call stack: factorial(5) -> 5 * factorial(4)
#             factorial(4) -> 4 * factorial(3)  ... -> factorial(1) = 1
# Unwinds:    1 -> 2 -> 6 -> 24 -> 120
\end{lstlisting}

\begin{lstlisting}
import sys

# Default recursion limit
print(sys.getrecursionlimit())   # 1000

# Change it (use carefully)
sys.setrecursionlimit(3000)
\end{lstlisting}

\begin{dangerbox}
Each recursive call adds a frame to the call stack. Too many recursive calls -> \textbf{RecursionError: maximum recursion depth exceeded}. For large inputs, prefer an iterative solution.
\end{dangerbox}

% -----------------------------------------------------------------------------
\section{Exception Handling}
% -----------------------------------------------------------------------------

\subsection{Types of Errors}

\begin{center}
\begin{tabular}{lll}
\toprule
\textbf{Type} & \textbf{When} & \textbf{Example} \\
\midrule
Syntax Error & At parse time & \texttt{deff foo():} \\
Logical Error & Wrong output, no crash & Returning 2 instead of 1 in recursion base \\
Runtime Error / Exception & At runtime, unexpected input/state & Divide by zero \\
\bottomrule
\end{tabular}
\end{center}

\subsection{Exception Hierarchy}

\begin{verbatim}
BaseException
  +-- Exception
        +-- ArithmeticError
        |     +-- ZeroDivisionError
        +-- ValueError
        +-- TypeError
        +-- FileNotFoundError
        +-- IndexError
        +-- KeyError
        +-- RecursionError
        +-- ... (many more)
\end{verbatim}

\subsection{try / except / else / finally}

\begin{lstlisting}
a = int(input("Numerator: "))
b = int(input("Denominator: "))

try:
    result = a / b
    print(f"Result: {result}")
except ZeroDivisionError as e:
    print(f"Error: Cannot divide by zero. ({e})")
except ValueError as e:
    print(f"Error: Invalid input -- please enter numbers only. ({e})")
except Exception as e:
    print(f"Unexpected error: {e}")
else:
    print("Calculation successful!")  # runs only if NO exception
finally:
    print("End of execution.")   # ALWAYS runs -- clean up resources here
\end{lstlisting}

\begin{notebox}
\textbf{The \texttt{finally} block:} Always executes, whether or not an exception occurred. Perfect for closing files, database connections, or network sockets -- you must close resources even if your code crashes.
\end{notebox}

\subsection{Raising Exceptions Manually}

\begin{lstlisting}
def set_age(age):
    if age < 0 or age > 150:
        raise ValueError(f"Invalid age: {age}. Must be 0-150.")
    return age

try:
    set_age(-5)
except ValueError as e:
    print(e)   # Invalid age: -5. Must be 0-150.
\end{lstlisting}

% -----------------------------------------------------------------------------
\section{Concurrency: Threads \& Multiprocessing}
% -----------------------------------------------------------------------------

\subsection{Why Concurrency?}

Modern applications do many things at once: download files, play music, update UI, respond to user clicks. Without concurrency, each task would block the others.

\begin{itemize}
  \item \textbf{Thread:} Lightweight unit of execution. Multiple threads share the same process memory.
  \item \textbf{Process:} Heavier. Has its own memory space and Python interpreter (own GIL).
\end{itemize}

\subsection{Threads with \texttt{threading} Module}

\begin{lstlisting}
import threading
import time

def download(filename):
    print(f"Downloading {filename}...")
    time.sleep(0.5)   # simulate download time
    print(f"Downloaded {filename}!")

files = ["video.mp4", "image.png", "data.csv"]

# Serial (no threads)
start = time.time()
for f in files:
    download(f)
print(f"Serial: {time.time() - start:.2f}s")  # ~1.5s

# Parallel (with threads)
threads = []
start = time.time()
for f in files:
    t = threading.Thread(target=download, args=(f,))
    threads.append(t)
    t.start()           # LAUNCH each thread

for t in threads:
    t.join()            # WAIT for all threads to finish

print(f"Parallel: {time.time() - start:.2f}s")  # ~0.5s
\end{lstlisting}

\begin{lstlisting}
# Class-based approach (inheriting Thread)
class Hello(threading.Thread):
    def run(self):           # must be named 'run'
        for i in range(5):
            print(f"Hello {i+1}")

class Hi(threading.Thread):
    def run(self):
        for i in range(5):
            print(f"Hi {i+1}")

t1, t2 = Hello(), Hi()
t1.start()
t2.start()
t1.join()
t2.join()
print("bye")
\end{lstlisting}

\subsection{The GIL Problem}

Python's \textbf{Global Interpreter Lock (GIL)} allows only one thread to execute Python bytecode at a time. This means:

\begin{center}
\begin{tabular}{lll}
\toprule
\textbf{Task Type} & \textbf{Threads help?} & \textbf{Why} \\
\midrule
I/O-bound (file, network) & \textbf{Yes} & GIL released during I/O waits \\
CPU-bound (heavy math) & \textbf{No} & GIL prevents parallel CPU use \\
\bottomrule
\end{tabular}
\end{center}

\subsection{Multiprocessing for CPU-bound Tasks}

\begin{lstlisting}
from multiprocessing import Process
import time

def calculate(n1, n2):
    total = sum(i ** 2 for i in range(n1, n2))

num = 50_000_000
mid = num // 2

# Each process has its OWN GIL -- true parallelism
p1 = Process(target=calculate, args=(0, mid))
p2 = Process(target=calculate, args=(mid, num))

start = time.time()
p1.start(); p2.start()
p1.join();  p2.join()
print(f"Multiprocessing: {time.time() - start:.2f}s")
# Noticeably faster than serial for CPU-bound work
\end{lstlisting}

\begin{tipbox}
Use \code{threading} for I/O-bound tasks. Use \code{multiprocessing} for CPU-bound tasks. For truly async I/O at scale (web servers), look into \code{asyncio}.
\end{tipbox}

% -----------------------------------------------------------------------------
\section{Arrays (\texttt{array} module)}
% -----------------------------------------------------------------------------

\begin{lstlisting}
from array import array

# array(typecode, [values]) -- homogeneous typed array
arr = array('i', [10, 20, 30, 40, 50])   # 'i' = signed int

print(arr)         # array('i', [10, 20, 30, 40, 50])
print(list(arr))   # [10, 20, 30, 40, 50]  -- convert to list to print cleanly

# Operations
arr.append(60)             # add at end
arr.insert(1, 15)          # insert 15 at index 1
arr.remove(30)             # remove first occurrence of 30
val = arr.pop()            # pop last element
print(arr.typecode)        # 'i'
print(arr.buffer_info())   # (address, length)

# Copy an array efficiently (generator expression)
arr2 = array(arr.typecode, (n for n in arr))

# Iterate
for n in arr:
    print(n)
\end{lstlisting}

\begin{notebox}
\textbf{array vs list:} Arrays require all elements to be the same type -- faster for numeric operations and uses less memory. For complex array work (2D arrays, matrix ops), use NumPy instead.
\end{notebox}

% -----------------------------------------------------------------------------
\section{Quick Reference -- Patterns \& Idioms}
% -----------------------------------------------------------------------------

\subsection{List Comprehensions}

\begin{lstlisting}
# [expression for item in iterable if condition]
squares = [x**2 for x in range(10)]
evens   = [x for x in range(20) if x % 2 == 0]
matrix  = [[i*j for j in range(1,4)] for i in range(1,4)]
\end{lstlisting}

\subsection{Dictionary / Set Comprehensions}

\begin{lstlisting}
word_lengths = {word: len(word) for word in ["hello", "world", "python"]}
# {'hello': 5, 'world': 5, 'python': 6}

unique_squares = {x**2 for x in range(-5, 6)}
# {0, 1, 4, 9, 16, 25}
\end{lstlisting}

\subsection{enumerate \& zip}

\begin{lstlisting}
names = ["Naven", "Kiran", "Hush"]

for i, name in enumerate(names, start=1):
    print(f"{i}. {name}")   # 1. Naven, 2. Kiran, 3. Hush

scores = [95, 87, 72]
for name, score in zip(names, scores):
    print(f"{name}: {score}")
\end{lstlisting}

\subsection{Walrus Operator \texttt{:=} (Python 3.8+)}

\begin{lstlisting}
import re
data = "Hello 42 World"

if m := re.search(r'\d+', data):
    print(f"Found number: {m.group()}")   # Found number: 42
\end{lstlisting}

\subsection{Context Managers (\texttt{with})}

\begin{lstlisting}
# Automatically closes the file even if an exception occurs
with open("data.txt", "r") as f:
    content = f.read()
# file is closed here -- no need for explicit f.close()
\end{lstlisting}

\subsection{Unpacking \& Starred Assignment}

\begin{lstlisting}
first, *middle, last = [1, 2, 3, 4, 5]
print(first)    # 1
print(middle)   # [2, 3, 4]
print(last)     # 5
\end{lstlisting}

% -----------------------------------------------------------------------------
\section{Complete OOP Example -- Putting It All Together}
% -----------------------------------------------------------------------------

\begin{lstlisting}
from abc import ABC, abstractmethod

# ----------- ABSTRACT BASE CLASS (Interface) -----------
class Shape(ABC):
    @abstractmethod
    def area(self): pass

    @abstractmethod
    def perimeter(self): pass

    def __str__(self):
        return f"{self.__class__.__name__}: area={self.area():.2f}"

# ----------- CONCRETE CLASSES (Implementations) -----------
class Circle(Shape):
    PI = 3.14159

    def __init__(self, radius):
        self.radius = radius

    def area(self):
        return Circle.PI * self.radius ** 2

    def perimeter(self):
        return 2 * Circle.PI * self.radius

class Rectangle(Shape):
    def __init__(self, width, height):
        self.width  = width
        self.height = height

    def area(self):
        return self.width * self.height

    def perimeter(self):
        return 2 * (self.width + self.height)

# ----------- POLYMORPHISM IN ACTION -----------
shapes = [Circle(5), Rectangle(4, 6), Circle(3)]

for shape in shapes:
    print(shape)              # uses __str__
    print(f"  Perimeter: {shape.perimeter():.2f}")

# All shapes respond to area() and perimeter() -- duck typing!
total_area = sum(s.area() for s in shapes)
print(f"Total area: {total_area:.2f}")
\end{lstlisting}

% -----------------------------------------------------------------------------
\section{Mental Models -- Summary}
% -----------------------------------------------------------------------------

\begin{summarybox}
\begin{itemize}
  \item \textbf{Variables = labels, not boxes.} They point to objects. Objects have identity (\code{id}), type (\code{type}), value.
  \item \textbf{Everything is an object.} Integers, strings, functions, classes -- all objects.
  \item \textbf{Mutable vs Immutable.} Lists/dicts/sets are mutable. Strings/tuples/ints are immutable.
  \item \textbf{self = current instance.} Passed automatically. Write it, don't call it.
  \item \textbf{Inheritance chains.} MRO is left-to-right, depth-first. Use \code{super()} to call parent.
  \item \textbf{Duck typing.} Python doesn't care about the class, only whether the method exists.
  \item \textbf{Decorators = wrappers.} \code{@dec} above \code{def f} means \code{f = dec(f)} behind the scenes.
  \item \textbf{GIL.} Threads for I/O. Processes for CPU. \code{asyncio} for high-concurrency I/O.
  \item \textbf{Exceptions.} \code{try/except/finally} -- never let the app crash on predictable errors.
  \item \textbf{Modules = .py files. Packages = folders of modules.} \code{\_\_name\_\_ == '\_\_main\_\_'} guards execution.
\end{itemize}
\end{summarybox}

% --- Footer -------------------------------------------------------------------
\vfill
\begin{center}
{\small\color{gray}\itshape Python Notes -- Intermediate-level reference. Pair with code practice for best retention.}
\end{center}

\end{document}
